% =====================================================================
%  facilitator_quickref.tex
%
%  Two-page quick reference. Print double-sided on one sheet.
%  Page 1 = how to run it.  Page 2 = when to intervene.
%  Derived from facilitator_manual.tex; see that document for detail.
% =====================================================================

\documentclass[9pt,letterpaper]{extarticle}

\usepackage[letterpaper,top=0.5in,bottom=0.45in,left=0.55in,right=0.55in]{geometry}
\usepackage{mathptmx}
\usepackage[T1]{fontenc}
\usepackage[utf8]{inputenc}
\usepackage{microtype}
\usepackage{amssymb}
\usepackage{booktabs}
\usepackage{tabularx}
\usepackage{array}
\usepackage{xcolor}
\usepackage{enumitem}
\usepackage{multicol}
\usepackage[most]{tcolorbox}
\usepackage{titlesec}
\usepackage[hidelinks]{hyperref}

\definecolor{deep}{HTML}{1F3A5F}
\definecolor{warn}{HTML}{9C3B2E}
\definecolor{calm}{HTML}{2F6B4F}
\definecolor{soft}{HTML}{EEF1F5}

\pagestyle{empty}
\setlength{\parindent}{0pt}
\setlength{\columnsep}{14pt}
\renewcommand{\arraystretch}{1.12}
\newcolumntype{R}{>{\raggedright\arraybackslash}X}
\setlist{nosep,leftmargin=1.1em,itemsep=1pt}

\titleformat{\section}{\bfseries\normalsize\color{deep}}{}{0em}{}
\titlespacing*{\section}{0pt}{1.4ex}{0.5ex}

\newcommand{\say}[1]{\textit{``#1''}}
\newcommand{\hd}[1]{\vspace{0.5ex}{\bfseries\color{deep}#1}\vspace{0.3ex}\\}

\newtcolorbox{redbox}[1]{colback=warn!7,colframe=warn,fonttitle=\bfseries\small,
  title={#1},boxrule=0.7pt,arc=1.5pt,left=4pt,right=4pt,top=2.5pt,bottom=2.5pt,
  fontupper=\small}
\newtcolorbox{greenbox}[1]{colback=calm!7,colframe=calm,fonttitle=\bfseries\small,
  title={#1},boxrule=0.7pt,arc=1.5pt,left=4pt,right=4pt,top=2.5pt,bottom=2.5pt,
  fontupper=\small}
\newtcolorbox{graybox}[1]{colback=soft,colframe=deep!40,fonttitle=\bfseries\small,
  title={#1},boxrule=0.6pt,arc=1.5pt,left=4pt,right=4pt,top=2.5pt,bottom=2.5pt,
  fontupper=\small}

\begin{document}

% =====================================================================
% PAGE 1 -- HOW TO RUN IT
% =====================================================================
{\color{deep}\rule{\textwidth}{1.5pt}}\\[0.3em]
{\LARGE\bfseries\color{deep} Solutions Dimension --- Facilitator Quick Reference}\\[0.2em]
{\small p1 how to run it \quad$\cdot$\quad p2 when to intervene \quad$\cdot$\quad
p3 web version \quad$\cdot$\quad p4 scenario index \quad$\cdot$\quad Full detail in the Manual}\\[0.2em]
{\color{deep}\rule{\textwidth}{1.5pt}}

\vspace{0.6em}
\begin{graybox}{The one thing to understand}
\textbf{The disagreement is the lesson.} Students learning card definitions is a
side effect. The activity works when a team argues about which of two defensible
options is better and has to say why. \textbf{A session where everyone agreed
quickly did not work.} No security expertise is needed to facilitate.
\end{graybox}

\vspace{0.4em}
\begin{multicols}{2}

\section{Setup}
\begin{itemize}
\item Teams of \textbf{2--3}. Four is the maximum before someone stops talking.
\item Per team: original 4-dimension deck, Solutions deck, printed $3\times3$
risk matrix, \textbf{6 physical tokens}, record sheet.
\item Sort Solutions into piles before you start: Foundational (1--13),
Extended (14--18), Advanced (19--21), modifiers (22--28).
\item Use real tokens, not a tally. Spending must be visible and irreversible.
\item \textbf{Explain Phase 1 only, then start.} Explain each later phase when
teams reach it. Front-loading all three confuses everyone.
\end{itemize}

\section{Costs}
\small
\begin{tabularx}{\linewidth}{@{}l c X@{}}
\toprule
\textbf{Tier} & \textbf{Cost} & \textbf{Cards} \\
\midrule
Foundational & 1 & 1--13 \\
Extended & 2 & 14--18 \\
Advanced & 3 & 19--21 \\
Modifiers & 0 & 22--28 \\
\bottomrule
\end{tabularx}
\normalsize

\vspace{0.4em}
Budget \textbf{6 tokens}. Deal \textbf{8 control cards + 3 modifiers}. A good
default deal is 6 Foundational, 1 Extended, 1 Advanced. \textbf{Deal
deliberately, not randomly} --- if a team's threat calls for a card, include it.

\section{Timing (35--45 min total)}
\small
\begin{tabularx}{\linewidth}{@{}X r@{}}
\toprule
Setup and rules & 5--7 \\
Phase 1: threat and placement & 5--7 \\
Phase 2: budgeted defense & 8--10 \\
Phase 3: residual risk & 4--5 \\
Cross-team comparison & 5--8 \\
Debrief & 5--10 \\
\bottomrule
\end{tabularx}
\normalsize

\columnbreak

\section{Phase 1 --- Threat and Placement}
Draw one card from each of the four original dimensions. Build a scenario where
all four are true. Place on the matrix. Record the coordinates.

\vspace{0.3em}
\say{Build a story where all four of these are true at once. Every card has to
do some work in it.}

\vspace{0.3em}
\textbf{Severity rule.} They must justify severity by pointing at the
\textbf{Human Impact card} and naming who is harmed. Novices rate severity by
how sophisticated an attack \emph{sounds}; this corrects it without a lecture.

\section{Phase 2 --- Budgeted Defense}
Deal 8 Solutions + 3 modifiers. Give 6 tokens. For each purchase they say which
safeguard addresses which part of the threat.

\vspace{0.3em}
\say{Six tokens. Anything you don't spend is gone. For each card, tell me which
part of your scenario it touches.}

\vspace{0.3em}
\textbf{Enforce the arithmetic instantly.} A team that overspends unnoticed has
not played the game.

\section{Phase 3 --- Residual Risk}
Same threat, same board, re-placed given what they bought. They state what still
gets through.

\vspace{0.3em}
\say{Where does it sit now? And tell me what still gets through, because
something does.}

\vspace{0.3em}
\textbf{Then compare teams.} Don't skip it. The question is never who was right
--- it's why two reasonable groups bought different things.

\end{multicols}

\vspace{0.3em}
\begin{graybox}{Adjudicating modifier cards --- the one-sentence test}
Modifiers are free, so each requires a justification. If unsure, ask once:
\say{Can you say that again naming a specific person, system, or failure?}
Specific $\rightarrow$ accept. Still general after \textbf{one} prompt
$\rightarrow$ card returns to the deck. Don't prompt twice.
\vspace{0.3em}

\textbf{Accept:} \say{Staff will share the one MFA phone at the front desk.}
\quad \textbf{Reject:} \say{People are the weakest link.}
\vspace{0.2em}

\textbf{Accept:} \say{The exec director owns that decision, reviewed in twelve
months.} \quad \textbf{Reject:} \say{We accept the rest of the risks.}
\vspace{0.3em}

Also: \textbf{max 2 modifiers per round}; card 24 requires that they bought card
4 or 7; card 26 is played \emph{instead of} a purchase.
\end{graybox}

\vspace{0.3em}
\begin{multicols}{2}
\section{Debrief --- pick 3 or 4}
\begin{itemize}
\item Which purchase did you argue about most, and how did you settle it?
\item What did you want and couldn't afford? What would you tell the client
about that gap?
\item Did anything you bought only help \emph{after} the attack succeeded? Worth
the points?
\item Two teams bought differently for similar threats. Why?
\item Which card would you cut from this deck entirely?
\end{itemize}

\columnbreak

\section{Misconceptions worth naming}
\begin{itemize}
\item Controls \emph{reduce} risk; they don't remove it. (That's Phase 3.)
\item Expensive $\neq$ better. Cost tracks CIS Implementation Group, not value.
\item Awareness training is one anti-phishing layer; email filtering (12) is
cheaper and doesn't depend on a tired person at 4pm.
\item Backups limit consequence; the data is still disclosed.
\item Not everything must be fixed. Four treatments: avoid, mitigate, transfer
(21), accept (26).
\end{itemize}
\end{multicols}

\newpage
% =====================================================================
% PAGE 2 -- WHEN TO INTERVENE
% =====================================================================
{\color{deep}\rule{\textwidth}{1.5pt}}\\[0.3em]
{\LARGE\bfseries\color{deep} When to Intervene}\\[0.2em]
{\color{deep}\rule{\textwidth}{1.5pt}}

\vspace{0.5em}
\begin{graybox}{Your default is NON-intervention}
New facilitators intervene far too early. A team silent for 30 seconds is
thinking. A team arguing is learning. A team corrected by a peer has learned more
than one corrected by you. \textbf{Your standing job is small:} keep time,
enforce budget arithmetic, adjudicate modifiers, and handle the four red-flag
situations below. Everything else is optional and most of it is
counterproductive.
\end{graybox}

\vspace{0.4em}
\textbf{\color{deep}The escalation ladder} --- use the lightest rung that works, then stop.

\vspace{0.3em}
\small
\begin{tabularx}{\textwidth}{@{}l l X@{}}
\toprule
\textbf{Rung 1} & Wait & Count to 20 before deciding a silence is a stall. \\
\textbf{Rung 2} & Ask an open question & \say{What's the worst thing that happens in your scenario?} Redirects without supplying content. \\
\textbf{Rung 3} & Point at a card & Physically point at the Methods card and say nothing. Astonishingly effective. \\
\textbf{Rung 4} & State a rule & \say{Remember you have to say which part of the threat each purchase touches.} \\
\textbf{Rung 5} & Direct instruction & Only when time is nearly gone. Supplying an answer ends the learning for that team. \\
\bottomrule
\end{tabularx}
\normalsize

\vspace{0.7em}
\textbf{\color{deep}Triggers and thresholds}

\vspace{0.3em}
\small
\begin{tabularx}{\textwidth}{@{}p{4.4cm} p{1.9cm} X@{}}
\toprule
\textbf{Signal} & \textbf{Threshold} & \textbf{What to do} \\
\midrule
Total stall, nobody speaking & 60 seconds & Rung 2. \say{Read your Methods card out loud and tell me who it hurts.} Reading aloud restarts almost every stalled table. \\
One person doing all the talking & 2 minutes & Rung 2, aimed at someone else: \say{Which card would you defend that nobody has picked?} Don't name or silence the dominant speaker. \\
\textbf{Agreement in under 90 sec} & Immediate & Rung 2. Argue the other side: \say{Make the case for the card you just dismissed.} Fast consensus means the trade-off wasn't felt. \\
Jargon others don't share & First time & Rung 2, to the speaker: \say{Say that again for someone who hasn't taken a security course.} \\
Buying without reference to the threat & 2nd purchase & Rung 3: point at the Methods card. Then Rung 4: \say{Which part of the scenario does that touch?} \\
Overspending the budget & Immediate & Rung 4. Correct at once --- the constraint is the mechanic. \\
Modifier played with no justification & Immediate & Rung 4. One prompt, then the card returns to the deck. \\
Threat moved to Not Likely / Not Severe in Phase 3 & Immediate & Rung 2. \say{So this can't happen at all now? What would an attacker with more time do?} \\
\say{Just tell us the right answer} & After one deflection & Be honest: \say{There's a better and a worse answer, but not one right one. Tell me your reasoning and I'll tell you where it's weak.} \\
Drifting into how to perform the attack & First time & Rung 4. \say{We're on the defending side of the table.} No moralizing. \\
Finished 5+ minutes early & Immediate & Rung 2. \say{Your budget was just cut by two tokens. What goes?} \\
Running out of time & 3 min left & Rung 5. Make them commit to a portfolio even if unfinished --- incomplete is assessable, abandoned isn't. \\
\bottomrule
\end{tabularx}
\normalsize

\vspace{0.7em}
\begin{redbox}{Four situations --- act immediately, the ladder does not apply}
\textbf{1. Someone discloses a real vulnerability} at their employer, student
org, or family business. \textbf{Stop it at the table:} \say{Let's not work
through a real system in the open --- come find me after.} Follow up privately.
A live vulnerability discussed in a room of twenty creates risk for a third
party who didn't consent.
\vspace{0.4em}

\textbf{2. A scenario lands on someone's personal experience.} Human Impact
includes Physical and Emotional Wellbeing; Methods includes Manipulation or
Coercion. Stalking, doxxing, or intimate partner surveillance can hit hard.
Watch for someone going quiet or leaving. \textbf{Offer an exit without
requiring an explanation, addressed to the team, not the person:} \say{This team
can redraw the Impact card if you want a different scenario.} Never require
anyone to narrate aloud. Keep a redraw available every session as a standing,
unremarkable option.
\vspace{0.4em}

\textbf{3. Someone is being spoken over or dismissed} --- usually a non-major by
a technical student. Fix it structurally, not socially: \say{Let's hear each
person's top card before anyone argues.}
\vspace{0.4em}

\textbf{4. A factual error that will leave the room as a belief} --- e.g.\ that
encryption stops phishing, backups prevent a breach, or insurance removes the
duty to notify. Correct briefly, without ceremony, then hand it back. Leave all
other wrong answers to peers or the debrief.
\end{redbox}

\vspace{0.5em}
\begin{greenbox}{Leave these alone}
\begin{multicols}{2}
\begin{itemize}
\item Silence under 20 seconds. It's thinking.
\item Disagreement, including heated. It's the mechanism --- step in only if it
turns personal.
\item A defensible choice you wouldn't have made. Ask why; don't correct.
\item A wrong answer a student is about to catch. Wait three seconds.
\item A team ignoring a card you think is obviously right. Note it for debrief;
steering purchases hollows out the exercise.
\item Slightly unrealistic scenarios. Plausibility isn't the objective.
\item A team spending everything on one expensive card. Let Phase 3 teach it.
\end{itemize}
\end{multicols}
\end{greenbox}

\vspace{0.4em}
{\footnotesize\color{deep}\textbf{Self-check after your third session:} Did teams
argue? Did I go above Rung 2 more than twice per table? Did non-majors speak as
much as majors? Did I supply an answer a peer would have given within five
seconds?}


\newpage
% =====================================================================
% PAGE 3 -- RUNNING THE WEB VERSION
% =====================================================================
{\color{deep}\rule{\textwidth}{1.5pt}}\\[0.3em]
{\LARGE\bfseries\color{deep} Scenario Operations --- Web Version}\\[0.2em]
{\color{deep}\rule{\textwidth}{1.5pt}}

\vspace{0.5em}
\begin{graybox}{Everyone opens the same link}
Facilitator adds \texttt{\#facilitator} to the URL. Projector adds
\texttt{\#projector}. Players use the plain link. Add \texttt{?session=tue-am}
to run two classes at once --- sessions are completely separate.
\end{graybox}

\vspace{0.4em}
\begin{multicols}{2}

\section{Starting a session}
\begin{enumerate}[label=\arabic*.]
\item Open with \texttt{\#facilitator}, pick a code, \textbf{Claim the role}.
Write the code down --- it is how you get the role back if your browser closes.
\item A fresh session \textbf{loads Scenario 1 automatically}. To use a different
one, pick it from the dropdown and press \textbf{Load}.
\item Loading sets the organization, the read-aloud text, the budget, and moves
play to Phase 1. It does not touch teams.
\item Set \textbf{Teams} to about class size $\div$ 3, rounded up. Options are
6, 8, 10, 12, 16.
\item Read the scenario aloud. Do not paraphrase it --- teams comparing notes
later need to have heard the same thing.
\end{enumerate}

\section{During play}
\begin{itemize}
\item \textbf{Phase 2} --- press \textbf{Deal} once. Dealing again replaces every
team's hand.
\item \textbf{Timer} --- 5/8/10/15 min buttons; it shows on the projector and
turns amber under two minutes.
\item \textbf{Adjudication} --- modifier justifications queue here. \emph{Specific
enough} or \emph{Return to deck}. One prompt only.
\item \textbf{Phase 3} --- the scenario's own debrief questions appear in your
panel automatically.
\end{itemize}

\columnbreak

\section{The expert key}
Visible only to you, under the loaded scenario. It lists \textbf{strong},
\textbf{partial}, and \textbf{weak-here} controls.

\vspace{0.3em}
\textbf{Use it} to follow the discussion, to know which trade-off the scenario
was built around, and to score justifications afterwards.

\vspace{0.3em}
\textbf{Do not} read it out, hint at it, or steer purchases toward it. A team
that buys a ``partial'' card and defends it well has done better work than one
that buys a ``strong'' card without reasoning. \textbf{Never show it to players.}

\section{Ending}
\begin{enumerate}[label=\arabic*.]
\item \textbf{Export CSV} first, always. Neither reset exports for you.
\item \textbf{New round} --- clears threats, hands, placements, purchases; teams
keep their seats. Use between scenarios with the same class.
\item \textbf{End session and clear board} --- clears everything including team
membership, archiving a copy first. Use at the end of class.
\item \textbf{Campaign play: do neither.} The same organization returns next week
with drift applied.
\end{enumerate}

\end{multicols}

\vspace{0.2em}
\begin{redbox}{Three things that go wrong}
\textbf{Two people on one team both clicking.} The board is shared and the last
write wins, so one set of changes silently overwrites the other. The app names a
driver and warns everyone else. Agree who types before Phase 2.
\vspace{0.3em}

\textbf{Dealing twice.} It replaces hands, including cards a team already bought.
Deal once, at the start of Phase 2.
\vspace{0.3em}

\textbf{Clearing before exporting.} \emph{New round} does not archive. Export
first or the round is gone.
\end{redbox}

\vspace{0.4em}
\begin{greenbox}{If the app fails mid-class}
Keep playing on paper. Every mechanic works with the physical deck, a printed
$3\times3$ grid, and six coins --- the app is a convenience, not the activity.
Collect the record sheets instead of exporting.
\end{greenbox}

\newpage
% =====================================================================
% PAGE 4 -- SCENARIO INDEX
% =====================================================================
{\color{deep}\rule{\textwidth}{1.5pt}}\\[0.3em]
{\LARGE\bfseries\color{deep} Scenario Index}\\[0.2em]
{\small Full text, suggested draws and debrief questions: Appendix C of the
Facilitator's Manual. \quad $\bullet$ = may land on personal experience; offer
the redraw to the team \emph{before} reading.}\\[0.2em]
{\color{deep}\rule{\textwidth}{1.5pt}}

\vspace{0.6em}
\footnotesize
\begin{tabularx}{\textwidth}{@{}c p{3.5cm} c p{3.4cm} R@{}}
\toprule
\textbf{\#} & \textbf{Organization} & \textbf{Bgt} & \textbf{Strongest controls} & \textbf{The trap} \\
\midrule
1 & Harborview Food Bank \newline \textit{payroll redirection} & 6 &
Social Engineering Awareness; Deepfake Awareness; Email \& Browser &
Teams rate severity by how clever the attack was. It is severe because
\$18k is a month of cash. \\
2 & Cedar Ridge Health Clinic \newline \textit{ransomware} & 6 &
Backup \& Recovery; Incident Response; EDR &
The backup exists and is useless --- it is on the same server. Prevention buys
nothing here. \\
3 $\bullet$ & Wilder House shelter \newline \textit{stolen laptop} & 6 &
Encryption; Physical Security; Identity \& Access &
A login password feels like protection. It is not. \\
4 & Fairmount Land Trust \newline \textit{dormant account} & 5 &
Account Lifecycle; Identity \& Access &
Nobody attacked anyone. The right answer is the dullest card in the deck. \\
5 & Northbridge School District \newline \textit{vendor breach} & 10 &
Third-Party \& Contractual; Incident Response &
Half the deck is useless --- the district owns none of the affected systems. \\
6 & Millbrook Water Authority \newline \textit{debug mode live} & 10 &
Secure Configuration; Identity \& Access &
With ten points teams overspend. The fix costs one. \\
7 $\bullet$ & Eastgate Legal Aid \newline \textit{insider access} & 10 &
Audit Logging \& SIEM; Identity \& Access + Least Privilege; MDR &
Every preventive control fails. The access was authorized. \\
8 & Delridge Public Library \newline \textit{unknown device} & 6 &
Asset Inventory; Segmentation; Secure Configuration &
Segmentation is out of reach at six points. What do you tell the board? \\
9 & Ridgeway Arts Collective \newline \textit{unsupported software} & 6 &
Software Inventory; Deferred Investment; Backup &
Teams buy Software Updates. There are no patches --- that is the problem. \\
10 & Second Chance Rescue \newline \textit{shared mailbox} & 10 &
Password \& Authentication; Account Lifecycle; Incident Response &
The shared password exists because it is genuinely convenient. \\
\bottomrule
\end{tabularx}
\normalsize

\vspace{0.9em}
\begin{graybox}{Choosing one}
\textbf{First session:} 1, 2, 4 or 8. \quad \textbf{Once they know the loop:}
5, 6, 7, 10. \quad \textbf{Campaign play:} 9.
\vspace{0.3em}

Vary the teaching target between sessions. If nobody bought logging last time,
run 7, where detection is the only lever that works. If teams over-bought
expensive cards, run 6, where a one-point card wins.
\end{graybox}

\vspace{0.5em}
{\footnotesize\color{deep}\textbf{Never run a session against a real client's
systems.} All ten organizations are fictional. If a participant starts
describing a real vulnerability at a real place, stop it at the table and follow
up privately --- see page 2.}

\end{document}
